\section{Experiment 6: how far does the allocation signature reach?}
\label{sec:six}

Experiments 1--5 run on one scenario and six models from two organizations.
Two readings survive that evidence equally well. The allocation signature may be
a property of an agent acting under a verification budget, or it may be a
property of this scenario's structure and of how two laboratories post-train
their frontier line. Experiment 6 separates them.

It also tests the claim in this paper's title, which the earlier experiments
could not. With a budget of two lookups an agent can verify both the memory
behind its own plan and the memory the situation makes salient, so every result
so far is equally consistent with \emph{verification follows topical relevance,
and the plan's memory happens to be relevant too}. A budget of one forces the
choice, and the two accounts then predict opposite things.

Everything below was pre-registered in \path{workshop/PREREGISTRATION-EXP6.md}
before the corresponding calls, in eleven dated amendments that record what was
seen at each point, including the failures.

\subsection{Design}

Two new worlds, structurally isomorphic to Experiment 1 and audited to be so by
\path{workshop/scripts/exp6-audit.ts}: an on-call reliability agent at a
payments provider, and a sourcing agent at a hardware manufacturer. Each has six
inherited memories inside the same length band this paper uses throughout, five candidate
actions, exactly one action backed by two memories, one memory whose true
negative caveat the drifted arm has lost, and hidden source records carrying the
numbers.

Salience is manipulated in \SixDoses{} doses. \texttt{situationTempting}
differs from the base situation by exactly two substituted clauses, so those
clauses form a two-level factor and the four cells are derived mechanically
from the two registered strings; the audit asserts that dose 0 and dose AB
reproduce them byte for byte. The verification budget is 2 or 1.

\SixModels{} models across \SixOrgs{} organizations: Anthropic, OpenAI's
frontier line and its open-weight release, Google, xAI, DeepSeek, Alibaba, Meta,
Mistral, Moonshot, MiniMax and Tencent. \texttt{gpt-oss-120b} is deliberate: the
same organization as three models already tested, different weights and
post-training, so it separates the laboratory from the training regime.

\SixEpisodes{} episodes and \SixFailures{} schema violations in the analysed
set. A seventeenth model, \texttt{nvidia/nemotron-3.5-lightning}, is excluded
and reported in \S\ref{sec:sixlimits}; its \SixNotRunEpisodes{} episodes and
\SixNotRunFailures{} failures are held outside the analysed set and are not
counted in that zero. First pass only, as in Experiment 4.

\subsection{The on-path invariant is not a property of two laboratories}

At budget 2, the corrupted memory is verified in \SixOnPathTwoK/\SixOnPathTwoN{}
of the episodes whose first-pass intent it backs. \SixPerfectModels{} of the
\SixEvaluatedModels{} models are at exactly 100\% when pooled across cells, and
the lowest pooled rate is \SixWorstOnPathModel{} at \SixWorstOnPathPct\%.

H28 was registered on cells, not on pooled models: at or above 0.90 ``in every
model \emph{and every cell}'' with at least 20 intent-aligned episodes.
Evaluated that way, \SixHTwentyEightBelow{} of \SixHTwentyEightCells{} cells
falls below the floor --- \SixHTwentyEightWorstCell{}, at
\SixHTwentyEightWorstK/\SixHTwentyEightWorstN{} =
\SixHTwentyEightWorstPct\%. Pooling hides it, and pooling is not what was
registered, so both are given here.

The placement manipulation replicates and is larger than in Experiment 1. The
swap arm was run for \SixSwapModels{} models only, so the comparison is
restricted to those same \SixSwapModels{}: moving the same corruption onto the
memory backing the action the agent already intends takes verification from
\SixMatchedRelOffK/\SixMatchedRelOffN{} to
\SixMatchedRelOnK/\SixMatchedRelOnN{} in the reliability world and from
\SixMatchedProcOffK/\SixMatchedProcOffN{} to
\SixMatchedProcOnK/\SixMatchedProcOnN{} in procurement, against Experiment 1's
92/150 to 75/75.

\subsection{The off-path rate is close to a step function}

What moves with salience is the other half of the budget. H22 was registered as
a partial order, dose 0 below each single-clause dose and both below dose AB, and
it holds: \SixOffRelZero\%, \SixOffRelA\%, \SixOffRelB\%, \SixOffRelAB\% in
reliability and \SixOffProcZero\%, \SixOffProcA\%, \SixOffProcB\%,
\SixOffProcAB\% in procurement. It is not monotone in dose \emph{index} --- the
two single-clause doses are not ordered against each other, and in reliability
the second is far below the first --- and the registered hypothesis never
claimed it was.

The rise is steep. Two hand-written situations were tried before the dose series
and neither landed in the intermediate regime: the first left the intent-aligned
cell empty, the second put both cells at ceiling. That is why the dose series
replaced them, and H21, the manipulation check registered for the second
attempt, is a \emph{failed} pre-registered hypothesis, reported as such.

Experiment 1's within-arm contrast does reproduce, at the location H24 fixed in
advance. \SixQualifyingDoses{} dose-by-world cells satisfy the admissibility
rule, and H16 is supported in \SixHSixteenSupported{} of them, with on-path
minus off-path between \SixHSixteenMinGapPct{} and \SixHSixteenMaxGapPct{}
points. That rule is worth reading sceptically: because H23 puts the on-path
rate near 100\% everywhere, an admissibility band with an upper bound below
100\% makes H16 close to arithmetically safe wherever it applies, and moving
the bound from 80\% to 75\% would empty the qualifying set entirely. The
contrast reproduces; the rule that located it was not a demanding one.

\subsection{Budget 1: the invariant is a property of having a spare credit}

With two credits, \texttt{intent\_share} is exactly 0.5 in the large majority of
episodes: one credit goes to the agent's own plan and one goes to whatever else
the situation makes salient. The two credits do different jobs, which is why a
budget of two cannot separate intent from relevance.

With one credit the agent must choose. Measured now on \emph{any} memory backing
the stated intent rather than only on the corrupted one --- a different quantity
from the \SixOnPathTwoK/\SixOnPathTwoN{} above, and the one a single credit can
speak to --- it goes to the plan in \SixPlanOneK/\SixPlanOneN{} of episodes,
against \SixPlanTwoK/\SixPlanTwoN{} at budget 2. \emph{The memory behind the agent's own
plan is always checked} is therefore a statement about having a spare credit,
not about the allocation policy, and this paper's earlier phrasing is corrected
accordingly.

H25 was registered as a statement about cells: at least 80\% ``at every dose in
both worlds''. \SixHTwentyFiveBelow{} of \SixHTwentyFiveCells{} dose-by-world
cells fall below that floor, \SixHTwentyFiveLowCells{}. Read per model instead,
\SixBelowFloor{} of \SixModels{} is below it, \SixBelowFloorModel{}, at
\SixOpusPlanPct\%, against \SixOtherPlanLoPct\% to \SixOtherPlanHiPct\% for the
rest. The registered unit is the cell, and on that unit H25 fails in two of
eight; the per-model view locates where the failure concentrates and is not the
registered test.

\textbf{And the direction fails too, in one cell.} H25's falsification clause
named the condition: if the credit goes more often to the topically salient
off-path memory than to the agent's plan at any dose, the framing is wrong at
budget 1 and is to be reported as wrong here rather than in a limitation. On the
registered subset --- episodes where the salient memory does not back the stated
intent, with both quantities on that same denominator --- it does, at
\SixSalienceWonCell{}: the plan takes \SixSalienceWonPlanPct\% of
\SixSalienceWonN{} choices and the salient memory \SixSalienceWonSalPct\%.
It is \SixSalienceCellsWon{} of \SixHTwentyFiveCells{} cells, and pooled over
all of them the plan still leads \SixChoicePlanPct\% to
\SixChoiceSalientPct\%; but the cell where it reverses is the one where the
alternative has been made maximally live, which is the condition the title claim
is about. \SixBelowFloorModel{}, measured on that same common denominator, is
\SixOpusPlanMatchedPct\% plan against \SixOpusSalientMatchedPct\% salient.

An earlier draft of this section reported ``in 0 of 16 models does salience
win''. That figure compared a restricted numerator against an unrestricted
denominator and is withdrawn.

\subsection{What this does and does not establish}\label{sec:sixlimits}

It establishes that the allocation signature is not an artefact of one scenario
or of two organizations' post-training, that its ceiling is invariant to how
salient the alternative is made while its floor is not, and that under the one
budget that forces a choice, intent beats topical relevance in every model
tested.

It does not establish that the effect is exceptionless, and under one condition
it is falsified. \SixShortOfPerfect{} models are short of 100\% on-path at
budget 2 and one cell of \SixHTwentyEightCells{} is below the registered floor;
H25 fails in \SixHTwentyFiveBelow{} of \SixHTwentyFiveCells{} cells; and at
\SixSalienceWonCell{} topical salience beats the plan outright. It also does not
evaluate a remedy: the verification scheduler this paper argues for is still a
design implication and not a result.

Two comparisons deserve a caveat the reader should not have to find. The
placement contrast moves the corruption, but the swap arm also puts one visible
caveat on screen and lengthens that memory's body well beyond the band the other
five sit in, so it is not a one-variable manipulation in the way Experiment 4's
arms are. And the dose series conditions on \texttt{intent\_aligned}, which the
dose itself moves, so those cells share the post-treatment selection problem
\S\ref{sec:qualifier} raises about its own stratified analysis.

Two properties of the measurement are worth stating, with their limits.
\path{workshop/analysis/exp6_validate_stored.py} re-applies the runner's current
schema validator to every stored episode, including those written before that
validator existed, and none fails, so the frontier and open-model sets are
comparable on the schema. And of \SixCredits{} credits in the drifted arm,
\SixAmbiguous{} could be attributed to more than one memory. That second check
is weak by construction --- a credit names one id, so ambiguity would require a
model to emit two inside one string --- and it says nothing about whether a
credit was matched to the memory the model meant. The same audit reports
\SixUnresolvable{} credits that match no memory at all, which are scored as
spent and not on target.
